\documentclass[11pt]{article}
\usepackage[margin=1in]{geometry}
\usepackage{amsmath,amssymb}
\usepackage{bm}
\usepackage{booktabs}
\usepackage[colorlinks=true,allcolors=blue]{hyperref}
\usepackage{parskip}

\newcommand{\rk}{r_{\kappa}}
\newcommand{\Rev}{\mathcal{R}}
\newcommand{\logR}{\log \Rev}

\newenvironment{point}{\begin{quote}\itshape}{\end{quote}}
\newcommand{\reply}{\par\noindent\textbf{Response.}\ }

\title{\vspace{-2em}Response to reviewers\\[4pt]
\large Ms.\ Ref.\ No.: PHYSA-262689\\
\large ``The Ramsey community number as a renormalization-group crossing''\\
\large \textit{Physica A: Statistical Mechanics and its Applications}}
\author{Alexei Vazquez}
\date{}

\begin{document}
\maketitle

\section*{Summary}

I thank the Editor and the three reviewers for a careful and genuinely useful
reading. The manuscript has been substantially revised. Reviewer~1's criticism
of the presentation was well founded and the paper has been rewritten
throughout for narrative clarity: a plain-language abstract, a completely
rewritten introduction that defines every quantity before it is used and closes
with a roadmap of all sections, and an opening paragraph in every section
stating what that section is for. All specific points are addressed
individually below.

More importantly, three of the reviewers' questions asked for evidence that the
manuscript had asserted rather than demonstrated. In checking each of them I
found that the assertion was in fact \emph{wrong}, and the revision corrects it.
I want to flag these three corrections at the top, because they are the
substantive changes in this revision and they arose directly from the reports:

\begin{enumerate}
\item \textbf{Plain versus degree-corrected detection}
      (Reviewer~2, point~3). The original claimed that degree correction
      advances detection by two generations, from $\rk=684$ to $\rk=44$. Asked
      to confirm that the two comparisons use the same prior convention, I
      found that they did not. The plain evidence carried a node-labelling
      prior that the degree-corrected evidence did not. Once both are
      conditioned on the same prescribed partition, the plain and
      degree-corrected evidences agree to $O(\ln m)$ and give
      \emph{identical} $\rk$; the entire reported gap is the labelling prior,
      worth $-n\ln K$. This also explains, and makes universal, the slope ratio
      that the original paper had declared non-universal: it is $(1-n/m)^{-1}$,
      giving $3$ on the diamond and $2$ on a lattice with $n/m\to1/2$. New
      Sec.~4.2, new Table~2, revised Table~4, revised Fig.~3(a).

\item \textbf{The optimal number of communities} (Reviewer~1, point~9). Asked
      what the slope of the fit in the right-hand panel of the hierarchy figure
      is, I measured it: it is $0.250$, not $0.5$. The correct scaling is
      $K_{\rm opt}=\sqrt{\gamma}\,m^{1/4}\sim n^{1/4}$, not $\sqrt{n}$; the
      $\sqrt{\cdot}$ in the original was applied to $b^t=\sqrt{m}$ and then
      mistakenly identified with $n$. The invariant statement---optimal depth
      at half the depth of the lattice, $d_\star=t/2$---is unchanged. Revised
      Sec.~10, Eq.~(30), revised Fig.~5(b) with the fitted slope shown.

\item \textbf{The order of the finite-temperature transition} (Reviewer~1,
      point~8). Asked for evidence that the transition is continuous at $q=2$
      and first order for larger $q$, I computed two exact diagnostics and
      found neither supports it. The $10$--$90\%$ width of the step shrinks as
      $\lambda_T^{-t}$ for \emph{every} $K$ including $K=2$---ordinary
      finite-size scaling at a continuous transition---and the energy density
      is continuous through $T_c$ for every $K$, i.e.\ there is no latent heat.
      The transition is continuous throughout; the apparent jump was a
      finite-size effect. In the course of this I also obtained the exact
      large-$K$ asymptote $k_BT_c/J\to3/(2\ln K)$, replacing the previously
      fitted $c\approx0.70$ by the exact $2/3$. Rewritten Sec.~11, new Table~6,
      rebuilt Fig.~6.
\end{enumerate}

A fourth correction, prompted by Reviewer~2's point~2, concerns the
identification of the ordered partition: the state that orders below
$T_c^{\rm FM}$ is \emph{not} the bundle cut of Sec.~2, as the original text
said in passing, but the \emph{staggered} cut that gives one hub to each block.
The two sever the same $2^t$ seam edges and differ in energy by exactly
$\gamma$; this is now derived, verified by exhaustive minimisation, and drawn
(new Sec.~9.1, Eq.~(26), new Fig.~2(b)).

Every number in the revision is reproduced by the accompanying scripts, and a
new Appendix~A documents the three exact constructions in enough detail to
reimplement them. The revised manuscript is $33$ pages; a list of all changes
is given at the end of this letter.

\section*{Reviewer \#1}

\begin{point}
The authors describe a known merit of figure, the Ramsey community number from
the renormalization group point of view. The work is promising and potentially
of significant interest. However, it is very challenging to read it, even for
experts in the field. The presentation is poorly structured and narrated,
making the manuscript unnecessarily difficult to follow.
\end{point}

\reply I accept this criticism. The paper has been rewritten with the
presentation as the priority: every symbol is now defined in words before it
appears in an equation, every section opens by stating its purpose, and the
introduction ends with a roadmap. The specific points are addressed below.

\subsection*{1. The abstract}

\begin{point}
Already the Abstract is written as a far too difficult and technical text
without a proper narrative guideline. Such undefined parameters are used which
are known only by the very experts of the field. I would suggest to replace the
present version of the Abstract with a simplified version of the 3rd paragraph
of the Introduction.
\end{point}

\reply The abstract has been rewritten as suggested, in the narrative style of
the introduction's overview paragraph. It now opens with the question in plain
language (``a network has to be large enough before splitting it into
communities is a better description of it than not splitting it at all''),
defines $\rk$ verbally, and explains the two eigenvalues in words---``a volume
eigenvalue $bs$ that multiplies the bulk and a smaller seam eigenvalue $b$ that
multiplies the boundary between communities, so the boundary becomes relatively
thinner at every step''---rather than presenting a spectrum. Every symbol that
survives ($b$, $s$, $K$, $q$, $m$, $n$, $\gamma$) is defined where it appears.

\subsection*{2A. The fourth sentence of the Introduction: the two models, $\Rev$ and $P$}

\begin{point}
In the fourth sentence, the text presents a highly technical statement that
ends with the definition of two quantities, R and P. However, neither the two
models nor those two physical quantities are properly explained.
\end{point}

\reply The introduction now devotes a full paragraph to this before any formula
appears. The two models are described explicitly---``the first is a stochastic
block model: the nodes are divided into blocks, and the probability of an edge
depends only on which blocks its endpoints belong to, so the model can express
the statement `these groups exist'. The second is a null model with a single
block, in which every node is statistically equivalent to every other, so it
cannot express that statement at all.'' The word \emph{evidence} is then
defined (``each model assigns a probability to the observed network once its
parameters are integrated out; these two numbers are the evidences''), $\Rev$
is introduced as their ratio in Eq.~(1) with an explicit reading, and
$P(\text{split})$ is separated into its own Eq.~(2) with the decision rule
spelled out, including its equivalent form $\logR\ge\ln[q/(1-q)]$.

\subsection*{2B. What is ``the split''}

\begin{point}
The 5th sentence is similarly very puzzling. First of all, what is the split
mentioned in this sentence.
\end{point}

\reply The term is now defined explicitly at first use: ``Throughout, we call
the block-model hypothesis `the split'.'' It denotes the hypothesis that the
network is better described by a partition into blocks than by none; it is not
a particular partition being constructed, and no search over partitions is
performed anywhere in the paper (this is also now stated in the scope paragraph
of the introduction).

\subsection*{2C. The content of Eq.~(2)}

\begin{point}
The content, the definition of Eq.~(2) is also unclear after it is based on the
earlier mentioned unclear text.
\end{point}

\reply Equation~(2) of the original is now Eq.~(3), and it is preceded by a
verbal statement of exactly what it means: ``Detection fails on the small
members of the family and succeeds on the large ones, and the Ramsey community
number is the smallest size at which it succeeds.'' Since the two ingredients
it rests on ($\Rev$ and the decision rule) are now defined beforehand, the
definition should read cleanly. I have also added the reason for the name---the
analogy with Ramsey theory, where a structure is forced to appear once the
system is large enough---immediately after it.

\subsection*{3. Summary of the later sections}

\begin{point}
Some later sections are mentioned at the end of the Introduction, but it is
advisable to describe and briefly summarize all the later sections at the end
of the Introduction.
\end{point}

\reply The introduction now ends with a paragraph summarising every section in
turn, from Sec.~2 through the Conclusions and the Appendix, one clause each.

\subsection*{4A, 4B. The aim of each section}

\begin{point}
Similar statements can be repeated for Section 2 and 3: the narration is
missing from the text. What is the aim of Section 2? This must be briefly and
clearly mentioned at the beginning of Section 2 \& 3. All the sections must be
briefly modified in this sense.
\end{point}

\reply Every section now opens with a sentence stating its aim. For example,
Sec.~2: ``The aim of this section is to define the lattice, to define the
bipartition whose detectability the rest of the paper analyses, and to give in
closed form the handful of counts on which both models depend.'' Sec.~3: ``This
section establishes the structural fact on which everything else rests: growth
acts on the block statistics as an exact linear map, and the two eigenvalues of
that map separate bulk growth from the growth of the community boundary.'' The
same has been done for Secs.~4--12 and for the subsections of Secs.~6 and~11.

\subsection*{5. The $\lambda$ parameters in Eq.~(5)}

\begin{point}
In section 2, it is unclear what are those lambda parameters on the RHS in
Eq.~(5).
\end{point}

\reply This was poor presentation: the matrix and its eigenvalues were placed
side by side in one display without saying that the latter were the eigenvalues
of the former. They are now separated. Equation~(6) of the revision gives $M$
alone; the eigenvalues follow in Eq.~(7) after the sentence ``Because $M$ is
triangular, its eigenvalues can be read off the diagonal'', with the explicit
statement that ``they are eigenvalues of $M$ in the ordinary sense: any
component of $\bm v$ along the corresponding eigenvector is multiplied by that
factor at every generation'', followed by what each one physically controls
($\lambda_{\rm vol}$ the edge count and degree sums, $\lambda_{\rm seam}$ the
community boundary).

\subsection*{6. Numerical evidence that the last column of Table~2 tends to $\ln 2$}

\begin{point}
In Table 2, there is no numerical evidence presented that the last column
really gives the value ln2 in the large t limit.
\end{point}

\reply Agreed---the original table stopped at $t=10$, where the density is
$0.685$, which is not evidence for $0.693$. The revision devotes a separate
table (Table~3) to the convergence, carried to $t=20$ ($m=1.1\times10^{12}$
edges), and reports not only the density but the deficit
$\delta_t^{\rm exact}=1-\logR_{\rm dc}/(m\ln2)$ against the analytic prediction
$\delta_t$ derived in Sec.~7.4:

\begin{center}
\begin{tabular}{rrrrr}
\toprule
$t$ & $\logR_{\rm dc}/m$ & $\delta_t^{\rm exact}$ & $\delta_t$ (predicted) & ratio\\
\midrule
$10$ & $0.685387$ & $0.011196$ & $0.011175$ & $1.0019$\\
$14$ & $0.692494$ & $0.000943$ & $0.000943$ & $1.0001$\\
$20$ & $0.693133$ & $0.000020$ & $0.000020$ & $1.0000$\\
\bottomrule
\end{tabular}
\end{center}

against $\ln2=0.693147$. So the limit is not merely approached numerically; the
rate of approach is predicted analytically by the seam mode of the RG map and
confirmed to four significant figures. All entries are computed in $80$-digit
arithmetic from the closed forms.

\subsection*{7. The ``detection problem''}

\begin{point}
Below Fig.~3, a detection problem is mentioned, but it is not fully clear what
is meant by this term.
\end{point}

\reply The term is now defined at that point: ``By the \emph{detection problem}
we mean, throughout, the comparison of the two models of Eqs.~(1)--(2) on this
family of lattices, regarded as a statistical-mechanics problem in its own
right: $\logR$ plays the role of a free energy, the statistics $\bm v$ the role
of couplings, and $M$ the role of the RG recursion. In that dictionary Eq.~(15)
is its fixed-point free-energy density.'' The same dictionary is then restated
compactly in Sec.~8.

\subsection*{8. Continuous versus jump in Fig.~6 (left)}

\begin{point}
On the LHS of Fig.~6, it is not evidenced that for q=2 the data show a
continuous transition and for larger q the data show a jump. The sharpening
mentioned in the label of the figure is not visible.
\end{point}

\reply The reviewer is right, and pursuing this showed the claim itself to be
wrong. I computed two exact diagnostics.

\emph{(i) Width.} The $10$--$90\%$ width of the step, in units of the critical
coupling, shrinks geometrically with generation at a rate that coincides with
the thermal eigenvalue $\lambda_T=f'(r_\star)$ of the exact Potts recursion:

\begin{center}
\begin{tabular}{rrrrr}
\toprule
$K$ & $2$ & $4$ & $8$ & $16$\\
\midrule
measured width ratio per generation & $1.677$ & $1.852$ & $2.045$ & $2.250$\\
$\lambda_T$ from the fixed point     & $1.679$ & $1.853$ & $2.045$ & $2.250$\\
\bottomrule
\end{tabular}
\end{center}

This is ordinary finite-size scaling at a \emph{continuous} transition,
$\Delta K/K_c\sim L^{-1/\nu}$, and it happens at $K=2$---which is known to be
continuous---exactly as at $K=16$. A step therefore forms for every $K$, and
its formation is not evidence of a first-order transition.

\emph{(ii) Latent heat.} I computed the energy density (the fraction of
same-colour bonds) by propagating derivatives analytically through the
recursion, at $t=18$, and measured $\epsilon(K_c^+)-\epsilon(K_c^-)$ over
windows $\pm\delta$ with $\delta$ shrinking. It tends to zero for every $K$
($0.010,0.001,0.0001$ at $K=2$; $0.118,0.026,0.005$ at $K=16$, for
$\delta=10^{-2},10^{-3},10^{-4}$): there is no latent heat at any $K$.

Section~11 has therefore been rewritten. The transition is continuous for every
$K$ on this lattice, with a $K$-dependent thermal eigenvalue and exponent
$\nu(K)$ falling from $1.338$ at $K=2$ towards $1/2$; the step is a finite-size
crossover. Figure~6 has been rebuilt to show exactly this: panel (a) the
steepening with size at fixed $K$, panel (b) the width against generation for
$K=2,4,8,16$ with the $\lambda_T^{-t}$ prediction drawn through the points,
panel (c) the ordering temperature. The new Table~6 gives $r_\star$,
$\lambda_T$, $\nu$ and $T_c$ across nine decades of $K$. The physical
conclusion of the section---that each level of the hierarchy is thermally rigid
below its own $T_c$ and that $T_c$ falls as the hierarchy deepens---is
unchanged and now better supported.

\subsection*{9. The slope of the fit in Fig.~5 (right)}

\begin{point}
On the RHS of Fig.~5, what is the slope of the fit on the red color data, and
what conclusion can be made from the result of the fitting?
\end{point}

\reply This question exposed an error. The measured slope is $0.250$, not the
$0.5$ implied by the $\sqrt{n}$ guide line drawn in the original figure. The
correct result is
\[
  K_{\rm opt}=2^{d_\star}=\sqrt{\gamma}\,2^{t/2}=\sqrt{\gamma\,b^{t}}
  =\sqrt{\gamma}\,m^{1/4}\simeq\sqrt{\gamma}\,(3n/2)^{1/4},
\]
i.e.\ $K_{\rm opt}\sim n^{1/4}$, with blocks of $\sim n^{3/4}$ nodes. The
original wrote $2^{t/2}\sim\sqrt{n}$, which conflates $b^t=\sqrt{m}$ with $n$;
on this lattice $n\simeq\tfrac23 m$, so $2^{t/2}=m^{1/4}$, a quarter power of
the size and not a square root. Section~10 has been corrected accordingly, and
Fig.~5(b) now draws $m^{1/4}=(3n/2)^{1/4}$ as the prediction, keeps $\sqrt{n}$
as a dotted contrast curve so the difference is visible, and prints the fitted
slope ($0.250$ over the five largest sizes) on the panel.

Two remarks on what conclusion follows. First, the optimisation result that
does \emph{not} depend on the lattice is the depth: $d_\star=\tfrac12(t+\log_2
\gamma)$, the optimal cut is at half the depth of the hierarchy. The exponent
relating $K_{\rm opt}$ to $n$ follows from that plus the lattice's own relation
between $b^t$ and $n$, and is therefore construction-specific; the revision now
says so explicitly rather than claiming a general $\sqrt{n}$ law. Second, the
qualitative conclusion is untouched: the optimal number of communities grows
without bound with the network, so the number of Potts states that best
describes the network is not a constant of the model.

\subsection*{10A. The two interpretations of $\rk$ in the Conclusion}

\begin{point}
At the beginning of the Conclusion the author mentions two interpretations of
the Ramsey community number and the fitted slope version is ruled out but in
the text earlier it is not shown or evidenced that this version is indeed
wrong.
\end{point}

\reply Correct: the claim was made only in the Conclusions. The revision adds
Sec.~6.3, ``Two readings of $\rk$'', which states both readings and quantifies
the difference, and Sec.~6.2, which delimits the validity of the asymptotic
formula with a dedicated table (Table~5) comparing the estimated and exact
crossing generation across the $(b,s)$ family and over six thresholds. The
substance is:

\begin{itemize}
\item the slope reading is asymptotically correct, and in fact this paper
      removes the fit from it: Eq.~(15) \emph{derives} the slope, $\sigma=\ln
      K$ per edge, with no free parameters;
\item but as a way of \emph{locating} $\rk$ it fails in exactly the regime used
      in practice. At $q=0.99$ it places the crossing of the diamond at
      $t_\star=2$ ($n=12$) where the exact rule gives $t_\star=3$ ($n=44$); for
      $(b,s)=(3,2)$ it is wrong by two generations. It becomes reliable only
      for $L=\ln[q/(1-q)]\gtrsim10^5$, i.e.\ certainties $1-q\lesssim
      e^{-10^5}$;
\item the reason is quantitative and now tabulated: at the generations where a
      practical threshold is crossed, the evidence density is short of $\ln K$
      by $63\%$ ($t=3$) and $37\%$ ($t=4$), so an asymptotic slope
      overestimates the accumulated evidence and underestimates $\rk$.
\end{itemize}

The Conclusions now refer to this evidence instead of asserting it, and are
worded as a statement about where the two readings part company rather than
about one being simply wrong.

\subsection*{10B. The last sentence of the Conclusions}

\begin{point}
On the one hand, the last sentence has a bit ugly starting. On the other hand,
the problem mentioned by the author is not a simply one question but most
likely a circle of phenomena as different versions of the disordered networks
exist. Therefore, it is advisable to improve also this sentence.
\end{point}

\reply Agreed on both counts. The sentence beginning ``And whether $M$
survives\ldots'' is gone. In its place the Conclusions now close with a
paragraph that treats disorder as the family of questions it is, distinguishing
three cases and saying what is expected of each: (i) bond disorder on an exact
hierarchical lattice, where the recursion survives but acts on distributions
rather than numbers, so that the seam eigenvalue becomes a typical rate---the
Migdal--Kadanoff scheme handles this and it is the natural first step; (ii)
structural disorder in the growth rule, where $b$ and $s$ are themselves random
and the eigenvalues become growth rates of a multiplicative process, raising
the question of whether the evidence density still self-averages to $\ln K$;
and (iii) genuinely random graphs with planted communities, where no exact
decimation exists and detectability is itself a sharp phase transition, so the
question is whether the crossing picture survives on average with $\ln K$
replaced by an annealed or quenched density. I state which of these I regard as
tractable extensions of the present method and which is the substantive open
problem.

\section*{Reviewer \#2}

\begin{point}
The manuscript contains an original and technically interesting exactly
solvable treatment of community detectability on hierarchical networks. The
principal combinatorial and Bayesian results appear valuable. I like the paper
and think it is very interesting.
\end{point}

\reply Thank you. The three concerns are addressed below; the third in
particular led to a substantive correction.

\subsection*{1. Exact crossing versus asymptotic estimate; finite-size validity; $q=0.5$}

\begin{point}
Please distinguish consistently between the exact crossing definition in
Eq.~(12) and the asymptotic estimate in Eqs.~(13)-(14). The latter follows from
an asymptotic evidence density and should not be presented as an unqualified
exact closed-form formula. Its finite-size validity and the special case q=0.5
should be clarified.
\end{point}

\reply Agreed, and the distinction now structures Sec.~6, which is split into
``The exact crossing'' (6.1), ``The asymptotic estimate and when it is exact''
(6.2), and ``Two readings of $\rk$'' (6.3).

\emph{Status of each equation.} Equation~(16) (the former Eq.~12) is the
definition and is exact: $\logR_{\rm dc}$ is known in closed form at every $t$,
so $t_\star$ is obtained by evaluating a known function on the integers, with
no asymptotics. Equation~(17) (the former Eq.~13) is now explicitly labelled an
\emph{estimate}: ``it uses the fixed-point density in place of the exact
evidence and therefore ignores the deficit $\delta_t$ and the $O(\ln m)$ Occam
terms''. Equation~(18) (the former Eq.~14, the closed form $\rk(b,s;q)$) is
stated as ``the asymptotic closed form \ldots subject to the validity
conditions of Sec.~6.2''.

\emph{Finite-size validity.} This is now quantified rather than asserted. The
new Table~5 compares the estimate with the exact crossing across six lattices
and six thresholds. The estimate is exact for every lattice tested once
$L\gtrsim10^5$; at the practical certainties $q=0.99$ and $q=0.999$ it misses
for most of them, by one generation and---when the crossing occurs in the first
few generations---by two. The largest threshold at which a miss is found is
$5.6\times10^2$ for $(2,2)$ and $4.2\times10^4$ for $(4,2)$. I have also
corrected the statement in Sec.~7.5: the original said the estimate ``exceeds
by exactly one generation otherwise'', which is true only in the asymptotic
regime; the text now says one generation asymptotically and up to two for early
crossings, consistent with Table~5.

\emph{The case $q=1/2$.} This is now treated explicitly as case (ii) of
Sec.~6.2. At $q=1/2$ the threshold is $L=0$, the argument of the logarithm in
Eq.~(17) vanishes and the estimate is undefined ($t_\star\to-\infty$). This is
not a removable singularity: at $q=1/2$ the crossing is decided not by the
extensive $\ln K$ term at all but by the sign change of the $O(\ln m)$ Occam
remainder during the transient---the moment at which the split first repays its
own complexity. The exact rule handles it without modification and gives
$t_\star=2$ for every member of the family examined ($\logR_{\rm dc}=-0.25$ at
$t=1$ and $+0.03$ at $t=2$ for the diamond), hence
$\rk=12,23,38,30,56,62$ for $(b,s)=(2,2),(3,2),(4,2),(2,3),(2,4),(3,3)$. The
text now states that Eq.~(17) should be used only for $q>1/2$, and
quantitatively only for $L\gg\ln K$.

\subsection*{2. Details of the finite-temperature recursion; how the RB term enters}

\begin{point}
Please provide sufficient details of the finite-temperature recursion for
reproducibility. Explain how the degree-weighted Reichardt-Bornholdt term is
incorporated into the recursion and how it leads to the different-community hub
order, since Eq.~(24) by itself describes the local ferromagnetic Potts
recursion.
\end{point}

\reply This was the sharpest question in the reports, because the honest answer
is that Eq.~(24) does \emph{not} contain the degree-weighted term, and the
original manuscript did not say so. The revision separates the two cases
cleanly and adds a full appendix.

\emph{(a) $K=2$: the penalty is included exactly.} The penalty is infinite
range but depends on the configuration only through the integer
$M_\kappa=\sum_i\kappa_i\sigma_i$. The ferromagnetic partition function is
therefore computed \emph{resolved by} the value of $M_\kappa$, and the penalty
applied afterwards, $Z=\sum_M G(M)\exp[-\tfrac{\beta\gamma}{4m}(M^2-\sum_i
\kappa_i^2)]$. $G(M)$ comes from the same decimation with the bond weight
promoted from a number to a polynomial in a formal variable tracking
$M_\kappa$: the two-terminal weight of a height-$h$ sub-diamond is a
$2\times2$ array whose entries are dictionaries $\{M\mapsto\text{weight}\}$,
series composition convolves them while the internal spin (whose final degree
$2^h$ is fixed at birth) contributes $z^{\pm2^h}$, parallel composition
convolves the two bundles, and the poles of degree $2^t$ are attached last and
kept resolved so that $\langle\sigma_A\sigma_B\rangle$ is read off directly.
Nothing is approximated. This is now Appendix~A.2, together with the validation
(brute-force enumeration over all $2^n$ configurations at $t\le3$, and
cross-validation between two independent implementations at $t=4$).

\emph{(b) Why the hubs end up in different communities.} The original stated
this as a numerical output and, worse, identified the ordered state with the
bundle cut of Sec.~2---which is not correct, since the bundle cut puts both
hubs in one block. The revision derives it. Two bipartitions sever exactly the
same $c_t=2^t$ seam edges: the bundle cut, block~1 $=$ bundle~0 $\cup\{A,B\}$,
with degree imbalance $M_\kappa=2^{t+1}$; and the \emph{staggered} cut,
block~1 $=$ bundle~0 $\cup\{A\}$, block~2 $=$ bundle~1 $\cup\{B\}$, which is
exactly degree balanced, $M_\kappa=0$, because each block then carries
$\kappa=m$. Their ferromagnetic energies are identical, so
\[
  \mathcal{H}_{\rm bundle}-\mathcal{H}_{\rm stag}
   =\frac{\gamma}{4m}\big(2^{\,t+1}\big)^2=\gamma
\]
exactly, independent of size. Hence any $\gamma>0$ selects the staggered cut
and the hubs sit in opposite communities---which is precisely what the exact
solution shows. Brute-force minimisation over all $2^{n-1}$ bipartitions at
$t=2$ confirms the staggered cut is the global ground state ($E=-9.5$ at
$\gamma=1$, against $-8.5$ for the bundle cut and $-1.5$ for a single
community), and the identity is verified by direct construction through $t=7$.
This is the new Sec.~9.1 with Eq.~(26); Fig.~2 now shows both cuts side by
side. I also note there that the preference is $O(1)$ rather than extensive, so
the detection analysis---insensitive to the reassignment of two vertices---
applies to either cut.

\emph{(c) $K>2$: what Eq.~(24) does and does not contain.} Section~11 now opens
with an explicit statement. At $K=2$, where both parts are treated exactly, the
ordering temperature coincides with the ferromagnetic $T_c^{\rm FM}$ for every
$\gamma>0$: the penalty selects \emph{which} colouring orders, the
ferromagnetic term sets \emph{when} it orders. For the depth-$d$ hierarchy we
therefore compute the rigidity of the colour field from the ferromagnetic
$K$-state Potts model, having established in Sec.~10 that the RB energy at
resolution $\gamma$ selects the depth-$d_\star$ partition as its ground state.
The text says in as many words: ``Equation~(31) is exactly that local
ferromagnetic recursion, and it does not contain the penalty; the penalty
enters the argument only through the ground-state selection of Secs.~9.1
and~10. We state this explicitly to avoid any suggestion that the
infinite-range term has been decimated---it has not, and it need not be, since
it does not move $T_c$.''

\emph{(d) Reproducibility.} Appendix~A now documents all three constructions:
A.1 the statistics and evidence (construction, $80$-digit arithmetic, why that
precision is needed, agreement with closed forms to $10^{-6}$); A.2 the $K=2$
RB solution as above; A.3 the $K$-state Potts hierarchy, including the reduced
variable $r=B/A$, the single map $r'=f(r)$ it produces, the location of the
unstable fixed point, and the analytic derivative propagation used for the
energy density. Each subsection names the script that implements it. I have
also made the order parameter unambiguous: it is the hub--hub colour rigidity
$\rho_K=[P(\sigma_A=\sigma_B)-1/K]/(1-1/K)$, defined in Eq.~(33). (In the
original, the figure plotted the same-colour excess while the text described it
as the different-colour excess; this is now consistent.)

\subsection*{3. Is the evidence conditional on the partition? Same prior convention?}

\begin{point}
Please state whether the Bayesian evidence is conditional on the prescribed
bundle partition or includes a partition prior, and confirm that the plain and
degree-corrected comparisons use the same prior convention.
\end{point}

\reply I checked, and I cannot confirm it: they did not. This is the most
consequential correction in the revision, and I am grateful for the question.

\emph{What the convention is.} Both evidences are conditional on the prescribed
bundle partition: the block assignment is fixed by the construction of the
lattice, no prior over partitions is included, and no search over partitions is
performed. This is now stated explicitly at the head of the new Sec.~4.2, with
the reason (the partition is not an inference target; the question asked is
whether \emph{this} partition describes the network better than none).

\emph{What was inconsistent.} The plain-model evidence, following the companion
papers, additionally carried a node-labelling prior obtained by integrating a
Beta--Bernoulli membership probability, contributing
\[
  \Lambda_t=\ln B(n_1{+}\alpha,n_2{+}\alpha)-\ln B(\alpha,n{+}\alpha)
           =-n\ln K+O(\ln n)
\]
(numerically $\Lambda_t/(-n\ln2)=0.999999$ at $t=12$). The degree-corrected
evidence did not carry it. Since $\Lambda_t$ is extensive in $n$ while the
evidence is extensive in $m$, and $n/m\to2/3$ here, this does not merely shift
the comparison---it rescales its density.

\emph{Consequences, all now in the paper.} Putting both models in both
conventions (new Table~2, revised Fig.~3(a)) gives:

\begin{itemize}
\item Conditional on the same partition, the plain and degree-corrected
      evidences agree to $O(\ln m)$: their difference is $-1.5$, $-11.5$,
      $-22.5$ at $t=4,8,12$ against evidences of order $10^2$, $10^4$, $10^7$.
      The reason is structural and now stated in Sec.~2: the bundle cut is
      balanced in both node number and degree,
      $\kappa_1/\kappa_2=1+2(1/s)^t\to1$, so degree correction has nothing to
      remove.
\item Consequently degree correction does \emph{not} move $\rk$ on this
      lattice: both models give $\rk=44$ in the conditional convention and
      $\rk=684$ with the labelling prior (revised Table~4). The claim
      ``degree correction advances detection by two generations'' has been
      removed from the abstract, introduction, Sec.~6 and the conclusions, and
      replaced by the correct statement.
\item In the labelled convention both densities converge to
      $(1-n/m)\ln K$, i.e.\ $\ln2/3=0.2310$ here against a measured slope
      $0.2302$. The ratio of the two densities is therefore $(1-n/m)^{-1}$:
      exactly $3$ for the diamond ($n/m\to2/3$) and $2$ for a lattice with
      $n/m\to1/2$ such as the pseudofractal web---which is the factor of two
      reported there. The original had declared that factor non-universal; it
      is universal, once one identifies what it is a property of, namely the
      node-to-edge ratio and the labelling prior, and not degree correction.
\end{itemize}

\emph{What is unaffected.} None of the renormalization-group results depend on
the choice: $\Lambda_t$ is a fixed function of $n$ that shifts the threshold
crossing without touching the map $M$, its eigenvalues, the flow, or the
fixed-point density. The paper states this and uses the conditional convention
throughout, reporting the other one alongside wherever a number depends on it
(Tables~2 and~4, Fig.~3(a)).

\section*{Reviewer \#3}

\begin{point}
This manuscript presents a theoretically elegant and technically solid
contribution\ldots I find the work original and valuable\ldots Overall, this is
a strong and well-executed theoretical paper. I recommend acceptance.
\end{point}

\reply Thank you for the careful summary of the argument and for the
recommendation. The three suggestions are all adopted.

\subsection*{1. The notation $q$}

\begin{point}
The author may wish to clarify the notation q, since it is used both for the
Bayesian detection threshold and for the number of Potts states in later
sections.
\end{point}

\reply Adopted, by resolving the clash rather than explaining it. Throughout
the revision, $q$ denotes the detection certainty and nothing else, and $K$
denotes the number of communities everywhere---as the number of blocks in the
evidence, as the value of the fixed-point density $\ln K$, and as the number of
Potts states in Secs.~10 and~11. This has the additional benefit of making the
continuity visible: the $K$ whose logarithm is the fixed-point evidence density
is literally the $K$ of the $K$-state Potts model that orders at $T_c(K)$. The
introduction now states the convention explicitly, and $q_{\rm opt}$ has become
$K_{\rm opt}$, $T_c(q)$ has become $T_c(K)$, and Eq.~(31) carries $(K-1)$ and
$(K-2)$.

\subsection*{2. A paragraph on scope}

\begin{point}
It would also be helpful to add a short explanatory paragraph emphasizing that
the work is not intended as a general-purpose algorithm, but as an exactly
solvable theoretical model that clarifies the physical meaning of
detectability.
\end{point}

\reply Adopted. The introduction now contains, before the roadmap: ``this is a
study of an exactly solvable model, not a detection algorithm. The partition
tested is prescribed by the construction of the lattice, and no search over
partitions is performed or needed; the purpose is to expose, in a setting where
every quantity is exact, what the detectability of communities means
physically.'' The same point is reinforced in Sec.~4.2, where the conditional
prior convention is justified precisely on the grounds that the partition is
not an inference target.

\subsection*{3. Expanded discussion of disordered and random networks}

\begin{point}
Finally, the discussion on possible extensions to disordered or random networks
could be slightly expanded, since this is likely to be of broad interest to the
community.
\end{point}

\reply Adopted, and expanded rather more than slightly---this also answers
Reviewer~1's point 10B. The Conclusions now close with a paragraph
distinguishing bond disorder (recursion intact, couplings become distributions,
tractable by Migdal--Kadanoff), structural disorder in the growth rule ($b$ and
$s$ random, eigenvalues become growth rates of a multiplicative process, with
self-averaging of the evidence density the question to settle), and random
graphs with planted communities (no exact decimation, detectability itself a
sharp transition, the question being whether the crossing picture survives on
average with $\ln K$ replaced by an annealed or quenched density), with a
statement of which I regard as tractable and which is the substantive open
problem.

\section*{List of changes to the manuscript}

\textbf{Presentation.}
\begin{enumerate}
\item Abstract rewritten in plain narrative form; every symbol defined at first
      use.
\item Introduction rewritten: separate paragraphs for the question, the two
      models and the evidence ratio, the definition of $\rk$, the results, the
      scope of the work, and a roadmap of all sections.
\item An opening paragraph stating the aim added to every section and to the
      subsections of Secs.~6 and~11.
\item Notation unified: $q$ is the detection certainty only; $K$ is the number
      of communities and of Potts states throughout.
\item Matrix $M$ [Eq.~(6)] and its eigenvalues [Eq.~(7)] separated, with the
      meaning of the eigenvalues stated.
\item ``The split'' and ``the detection problem'' defined at first use.
\end{enumerate}

\textbf{Corrections.}
\begin{enumerate}\setcounter{enumi}{6}
\item Prior conventions made consistent (new Sec.~4.2, new Table~2, revised
      Table~4, revised Fig.~3(a)): plain and degree-corrected evidences agree
      to $O(\ln m)$ and give the same $\rk$; the previously reported gap is the
      node-labelling prior $\Lambda_t\to-n\ln K$ [Eq.~(11)]. The claim that
      degree correction advances detection by two generations is removed.
\item Density ratio between conventions identified as $(1-n/m)^{-1}$
      [Eq.~(12)], explaining both the diamond's $3$ and the pseudofractal's
      $2$; the ``not universal'' claim is removed.
\item $K_{\rm opt}$ corrected from $\sqrt{n}$ to $\sqrt{\gamma}\,m^{1/4}\sim
      n^{1/4}$ [Eq.~(30)]; Fig.~5(b) redrawn with the correct guide, a contrast
      curve, and the fitted slope $0.250$ shown. The invariant statement
      $d_\star=t/2$ is retained and emphasised.
\item Order of the finite-temperature transition corrected: continuous for
      every $K$, with the step a finite-size crossover of width
      $\lambda_T^{-t}$ and no latent heat (rewritten Sec.~11, new Table~6,
      rebuilt Fig.~6).
\item The asymptote $k_BT_c/J\to3/(2\ln K)$ derived exactly [Eq.~(34)],
      replacing the fitted $c\approx0.70$.
\item Ordered partition identified as the staggered cut, not the bundle cut,
      with $\mathcal{H}_{\rm bundle}-\mathcal{H}_{\rm stag}=\gamma$ exactly
      (new Sec.~9.1, Eq.~(26), new Fig.~2(b)).
\item Order parameter of Sec.~11 made consistent between text and figure
      [$\rho_K$, Eq.~(33)].
\item Validity of the asymptotic crossing formula corrected: one generation
      asymptotically, up to two for early crossings (Sec.~7.5, Table~5).
\end{enumerate}

\textbf{New material.}
\begin{enumerate}\setcounter{enumi}{14}
\item Section~4.2 on prior conventions, with Table~2.
\item Table~3: convergence of the evidence density to $\ln K$ through $t=20$,
      against the analytic deficit.
\item Section~6.2 on the validity of the asymptotic estimate, including the
      $q=1/2$ case, with Table~5.
\item Section~6.3 on the two readings of $\rk$.
\item Section~9.1 on which partition is the ground state.
\item Table~6: the Potts critical fixed point across nine decades of $K$.
\item Appendix~A: computational methods for all three constructions.
\item Expanded closing discussion of disorder in the Conclusions.
\end{enumerate}

\textbf{Figures.} Fig.~2 now has two panels (bundle and staggered cuts);
Fig.~3(a) shows both models in both prior conventions; Fig.~5(b) is redrawn as
described; Fig.~6 is rebuilt with three panels. Figs.~1 and~4 are unchanged.

\textbf{Reproducibility.} Two new scripts accompany the revision:
\texttt{rb\_groundstate.py} (ground-state identification and the exact energy
gap $\gamma$) and \texttt{potts\_fixedpoint.py} (Potts fixed point, thermal
eigenvalue, transition width, energy density). \texttt{diamond\_rg.py} gains
the labelling-prior and conditional-plain evidence functions. Every table and
figure in the paper is regenerated by these scripts.

\end{document}
